\documentclass[doctype=article,language=hebrew]{omnilatex}

\title{תמיכה בשפה העברית ב-OmniLaTeX}
\subtitle{דוגמה לתבנית מימין לשמאל}
\author{דוד כהן}
\date{\today}

\begin{document}
\maketitle

\begin{abstract}
מסמך זה מציג דוגמה לתמיכה בשפה העברית במערכת OmniLaTeX. הוא כולל
פריסה מימין לשמאל עם גופנים עבריים והגדרות מתאימות לשפה.
\end{abstract}

\section{מבוא}

השפה העברית היא אחת השפות העתיקות ביותר שעדיין בשימוש. היא נכתבת
מימין לשמאל, דבר המצריך טיפול מיוחד במערכות הסדר טקסט. מערכת
OmniLaTeX מספקת תמיכה מלאה בכיוון זה באמצעות חבילת ה-bidi.

\section{מתמטיקה בטקסט עברי}

ניתן לכתוב משוואות מתמטיות בתוך טקסט עברי תוך שמירה על הכיוון
הנכון משמאל לימין. לדוגמה, משוואת אוילר:

\begin{equation}
  e^{i\pi} + 1 = 0
  \label{eq:euler}
\end{equation}

וכן סדרת טיילור:

\begin{equation}
  f(x) = \sum_{n=0}^{\infty} \frac{f^{(n)}(a)}{n!}(x - a)^{n}
  \label{eq:taylor}
\end{equation}

\section{רשימות}

\begin{itemize}
  \item פריט ראשון ברשימה
  \item פריט שני
  \item פריט שלישי עם הסבר נוסף
\end{itemize}

\begin{enumerate}
  \item שלב ראשון
  \item שלב שני
  \item שלב שלישי
\end{enumerate}

\section{טבלאות}

\begin{table}[htbp]
  \caption{השוואה בין מערכות הסדר טקסט}
  \begin{tabular}{llr}
    \textbf{מערכת} & \textbf{סוג} & \textbf{שנה} \\
    \hline
    TeX  & קוד פתוח & 1978 \\
    LaTeX & קוד פתוח & 1984 \\
    OmniLaTeX & קוד פתוח & 2024 \\
  \end{tabular}
\end{table}

\section{סיכום}

דוגמה זו מוכיחה שמערכת OmniLaTeX מסוגלת לטפל בטקסטים עבריים
ביעילות גבוהה, תוך שמירה על איכות הסדר ותמיכה במשוואות מתמטיות,
טבלאות ורשימות בכיוון הנכון.

\end{document}
